\scanpage{51}

La relation (125$'$) \unsure{montre} qu'on a

% slanted

\margin{\uline{NB} Il revient au $=$ de se donner l'op. de $Z'_n$ sur $\unsure{\Pi|\pi}$ (en faisant un $Z'_n$-torseur), ou l'opération $\lambda\mapsto T_\lambda$ de $Z'_n$ sur $S$, reliées par (125$'$) -- \unsure{chacune|une|l'une} détermine l'autre.}

\[
\begin{aligned}
 \{s,t\}\in A&\Longrightarrow
 \overbrace{\{T_\lambda s,T_\lambda t\}}^{\inserted{\{s',t'\}}}\in A\\
 &\text{et}\qquad \unsure{\tau|t}(\{s',t'\})=\lambda^2\,\tau(\{s,t\}),
\end{aligned}
\tag*{(127)}
\]

donc $T_\lambda$ est un \unsure{autom} de la structure définie par \deleted{\unsure{$(A,Z'_n),\tau$}} \inserted{$\{S,\Pi\text{ tors. sous }Z'_n,A\subset S\times S,\tau\}$,} ou ce qui revient au $=$, de la structure ($S$, $\Pi$ tors. sous $Z'_n$, $\Pi\xrightarrow{\Sigma}$ ens des structures triangulées sur $S$ comme ens de sommets), canoniquement associée à ladite structure.

\deleted{Revenons à la question de décrire $\wedge$ (118) en termes des autres structures, on notera pour ceci que l'on a, si}

% slanted

\margin{il faut y revenir ultérieurement !}

\[
 \deleted{\left\{\begin{smallmatrix}\{s,t,u\}\in\unsure{F|\mathcal F}=\coprod_{i\in\Pi}F_{\Sigma(i)}\\v\in S\end{smallmatrix}\right.}
 \quad\deleted{\Longrightarrow s\wedge v+t\wedge v+u\wedge v-}
 \tag*{\deleted{(128)}}
\]

J'aimerais comprendre comment \inserted{(éventuellement)} \uuline{une des structures} $\Sigma(i)$ \uuline{détermine toutes les autres}, et l'opération de $Z'_n$ sur l'ens de ces structures. \deleted{\unsure{L}}

Soit donc $P_0=(S,A_0,F_0)$ une structure de 2-polyèdre \unsure{connexe.|combin.} régulier, triangulé en l'occurrence, fini par dessus le marché !

Soit

\[
 G_0=\operatorname{Aut}(P_0)\lhook\joinrel\longrightarrow\unsure{\mathfrak S|\mathscr S}_S,
 \tag*{(128)}
\]

% archive sheet number: 50
